\section{Exact occupation differences at regular decision boundaries}
\label{sec:v78-seams}
The original guard contains no section-distance factor. Internal
section decisions therefore cannot be discarded by the guard estimate.
Their signed current has an exact algebraic structure which should
be retained before absolute values are taken. This structure does
not require pooling the exact occupation classes.

Work first on a compact regular endpoint domain $D$, away from
$g=0$, within one physical word and one fixed displacement. Choose
$D$ with a finite piecewise smooth stratification compatible with
the internal section interfaces; compact definable subdomains of
the physical graph provide such choices. The initial and terminal
section tests hold throughout $D$. Its cells have constant internal
section decisions, up to lower-dimensional strata. Write
\[
 a(q)=\frac{W^\varepsilon(q)}{4\pi c_*|B_m(q)|},\qquad
 Y=X_q,\qquad A(q)=\sum_{j=0}^{m-1}\one_{Y_R^*}(T_R^jq).
\]
The notation $T_R^jq$ means the physical trajectory represented by
these endpoints. The functions $a,Y,F,\vartheta$ have the same
traces on both sides of an internal section interface; only the
occupation changes. The locally Lipschitz guard is sufficient for
the divergence identity by approximation.

\begin{theorem}[An exact discrete divergence for internal decisions]
\label{thm:v78-decision-current}
For a variable $z$, form the polynomial-valued density
$(F,\vartheta)_\#(a z^A\one_D\,dq)$. The contribution of its
\emph{internal section interfaces} to its roof derivative is
\begin{equation}\label{eq:v78-section-factor}
                 \mathscr S_D(z)=(z-1)\mathscr F_D(z).
\end{equation}
Both sides are finite polynomial-valued signed measures on
$\R\times\mathbb S^1$. If
$\mathscr F_D(z)=\sum_n z^n\mathscr F_{D,n}$, then the contribution
to the original exact occupation $n$ is precisely
\begin{equation}\label{eq:v78-exact-difference}
             \mathscr S_{D,n}=\mathscr F_{D,n-1}-\mathscr F_{D,n}.
\end{equation}
The displacement, collision count, initial and terminal tests and
original source normalization are unchanged in every coefficient.
\end{theorem}
\begin{proof}
Orient an interface $\Sigma$ from the cell with occupation $a_-$
to the cell with occupation $a_+$, and let $N$ be that normal.
The negative outward flux from the first cell and that from the
second sum to
\[
 (F,\vartheta)_\#\left(
    a(Y\cdot N)(z^{a_+}-z^{a_-})\,d\mathcal H^1|_\Sigma\right).
\]
There is no absolute value on this signed expression. For nonnegative
integers $u,v$, define
\[
 P_{u,v}(z)=
 \begin{cases}
  \sum_{j=u}^{v-1}z^j,&v>u,\\
  0,&v=u,\\
  -\sum_{j=v}^{u-1}z^j,&v<u.
 \end{cases}
 \qquad z^v-z^u=(z-1)P_{u,v}(z).
\]
Sum the pushforwards with $P_{a_-,a_+}$ in place of the difference
to define $\mathscr F_D$. Changing the chosen normal exchanges
$a_-,a_+$ and changes both signs, so this definition is independent
of orientation. A coincident interface at which several visits
change is covered by the same identity; a net zero change contributes
zero. Intersections of distinct smooth interfaces of lower dimension
carry no surface flux. Compact regularity makes all these signed
measures finite. Taking the coefficient of $z^n$ proves
\eqref{eq:v78-exact-difference}.
\end{proof}

\begin{remark}[Outer tests and true singularities are separate]
\label{rem:v78-boundary-scope}
An initial or terminal section boundary need not have an admissible
source on both sides. Its flux is therefore not covered by
\eqref{eq:v78-section-factor}. Nor does the theorem identify a
first-hit switch or a grazing word transition with a regular section
interface. All such contributions remain in the full physical current
and in the exit term of \eqref{eq:v78-geometric-defect}.
On an exhaustion the polynomial identity holds at every compact
stage. Passage of the complete current to distributions follows
from source-mass convergence, as in the retained angular current
theorem. Convergence or a uniform bound for the isolated primitive
$\mathscr F_D$ does not follow from that fact and is not assumed.
\end{remark}

\subsection{The arithmetic phase of a discrete difference}
A factor $z-1$ is not small near a general spectral peak. The following
estimate records exactly what such a factor can and cannot supply.
Let $\xi=(u_1,u_2,b,v)$ be the four Fourier coordinates. In one
fixed peak chart centered at $a_j$ with occupation coordinate $v_j$,
suppose a scalar amplitude satisfies
\[
 |H_{m,j}(a_j+\delta)|\le C e^{-c m|\delta|^2}
\]
throughout that chart, with constants uniform in the permitted
parameters. Let $F_{m,j}(y)$ be its Fourier integral with factor
$(2\pi)^{-4}$, and let $S_{m,j}(y)$ be the same integral with
multiplier $e^{iv}-1$. Here $y=(k_1,k_2,t,n)$; no division by an
arithmetic reference is involved.

\begin{proposition}[Peak-aware coefficient difference]
\label{prop:v78-peak-difference}
Under the displayed amplitude bound,
\begin{equation}\label{eq:v78-peak-difference}
 S_{m,j}(y)=(e^{iv_j}-1)F_{m,j}(y)+O(m^{-5/2}),
\end{equation}
uniformly in all targets. Equivalently the covariant difference
$F_{m,j}(k,t,n-1)-e^{iv_j}F_{m,j}(k,t,n)$ is $O(m^{-5/2})$.
The ordinary difference is $O(m^{-5/2})$ only when its retained
leading term vanishes; $v_j=0$ modulo $2\pi$ is one such case.
\end{proposition}
\begin{proof}
On the chart,
\[
 |(e^{iv}-1)-(e^{iv_j}-1)|\le|v-v_j|\le|\delta|.
\]
The absolute value of the remaining Fourier integral is bounded by
$C\int_{\R^4}|\delta|e^{-cm|\delta|^2}\,d\delta$.
The change of variable $\delta=m^{-1/2}w$ gives $O(m^{-5/2})$.
Multiplication by $e^{iv}$ shifts the extracted occupation from $n$
to $n-1$, proving the second formulation. A moving center does not
change the estimate when the displayed constants are uniform.
\end{proof}

The hypothesis of this proposition is a scalar bound for the
\emph{surface-current primitive}, not just for the original volume
source. No such uniform bound for all physical interfaces is imported
from a theorem about bounded volume multipliers. Furthermore a
nonzero occupation peak retains the first term in
\eqref{eq:v78-peak-difference}, including at transition radii. This
prevents a spurious cancellation of nontrivial arithmetic classes.
The exact interface identity and its peak-aware estimate provide a
specific signed summation mechanism for the still required geometric
bound; they do not certify that bound or remove any Fourier peak.
